\section{Derivasi Leftmost dan Rightmost}
\label{sec:ch08-derivasi}

\subsection{Derivasi}

Derivasi adalah urutan penggantian nonterminal hingga diperoleh
string terminal. Pada setiap langkah, satu nonterminal diganti dengan
ruas kanan aturan produksi. Derivasi leftmost selalu mengganti
nonterminal paling kiri; derivasi rightmost mengganti nonterminal
paling kanan.

\begin{contoh}[Derivasi leftmost]
Untuk tata bahasa aritmetika dan string
\(\code{id} + \code{id} * \code{id}\):
\[
\begin{aligned}
E &\Rightarrow E + T\\
  &\Rightarrow T + T\\
  &\Rightarrow F + T\\
  &\Rightarrow \code{id} + T\\
  &\Rightarrow \code{id} + T * F\\
  &\Rightarrow \code{id} + F * F\\
  &\Rightarrow \code{id} + \code{id} * F\\
  &\Rightarrow \code{id} + \code{id} * \code{id}.
\end{aligned}
\]
\end{contoh}

\begin{contoh}[Derivasi rightmost]
Dengan tata bahasa yang sama:
\[
\begin{aligned}
E &\Rightarrow E + T\\
  &\Rightarrow E + T * F\\
  &\Rightarrow E + T * \code{id}\\
  &\Rightarrow E + F * \code{id}\\
  &\Rightarrow E + \code{id} * \code{id}\\
  &\Rightarrow T + \code{id} * \code{id}\\
  &\Rightarrow F + \code{id} * \code{id}\\
  &\Rightarrow \code{id} + \code{id} * \code{id}.
\end{aligned}
\]
\end{contoh}

\subsection{Bentuk Kalimat}

Setiap string yang muncul selama derivasi disebut bentuk kalimat
(\eng{sentential form}). Bentuk kalimat dapat memuat campuran
terminal dan nonterminal. String terminal yang dihasilkan disebut
kalimat (\eng{sentence}).

Sebagai contoh, pada derivasi di atas, \(E + T * F\) adalah bentuk
kalimat, sedangkan \(\code{id} + \code{id} * \code{id}\) adalah
kalimat.

\subsection{Pohon Derivasi}

Derivasi dapat direpresentasikan sebagai pohon. Akar adalah simbol
awal; setiap penerapan aturan menambahkan anak-anak pada simpul
nonterminal. Daun pohon, dibaca dari kiri ke kanan, membentuk string
yang diturunkan.

\begin{figure}[H]
\centering
\begin{forest}
parsetree
[E
  [E [T [F [\code{id}]]]]
  [+]
  [T
    [T [F [\code{id}]]]
    [*]
    [F [\code{id}]]
  ]
]
\end{forest}
\caption{Pohon urai untuk \(\code{id} + \code{id} * \code{id}\)}
\label{fig:ch08-parse-tree}
\end{figure}

\subsection{Derivasi dan Ambiguitas}

Satu string dapat memiliki lebih dari satu derivasi leftmost jika
tata bahasa ambigu. Ambiguitas tidak selalu terlihat dari string
akhir; ia terlihat dari pohon urai yang berbeda. Bab IX membahas
ambiguitas, dangling else, dan cara mengatasinya.

\begin{infobox}
Derivasi leftmost dan rightmost pada tata bahasa yang tidak ambigu
menghasilkan pohon urai yang sama. Perbedaan urutan penggantian tidak
mengubah struktur pohon.
\end{infobox}
